\section{Length cutoffs and open paths}\label{sec:cutoff-laws}

The analytic work is complete. We now use only the positive polygon estimates.

\begin{proof}[Proof of Corollary~\ref{cor:length-cutoffs}]
Use the functions $A$, $P$ and $M_2$ of Theorem~\ref{thm:polygon-tails}, whose bounds have now been proved. All sums below use the same positive polygon weights, and all exponent estimates hold as their argument tends to infinity. We prove

\[
F(n):=\sum_{L\ge n}w=n^{-3/2+o(1)},\qquad
 Q(n):=\sum_{L\ge n}Lw=n^{-1/2+o(1)}.
\]

First consider a length bin $n\le L<2n$ and split it at $D=n^{3/4}$. On its small-diameter part,

\[
\sum w\le n^{-2}M_2(n^{3/4})=n^{-3/2+o(1)}.
\]

On its large-diameter part,

\[
\sum w\le A(n^{3/4})=n^{-3/2+o(1)}.
\]

Multiplying these bounds by $2n$ gives a length-weighted bin bound $n^{-1/2+o(1)}$. Summing the bins $[2^j n,2^{j+1}n)$, $j\ge0$, gives

\[
F(n)\le n^{-3/2+o(1)},\qquad Q(n)\le n^{-1/2+o(1)}.
\]

To justify the infinite sum, first bound every $o(1)$ by an arbitrary fixed positive exponent loss, small enough that both geometric series converge, and then let that loss tend to zero.

For the reverse bounds fix $\delta>0$ and put $R=n^{3(1+\delta)/4}$. The length-weighted mass at diameter at least $R$ is

\[
P(R)=n^{-(1+\delta)/2+o(1)}.
\]

Its contribution from $L<n$ is at most

\[
nA(R)=n^{-1/2-3\delta/2+o(1)},
\]

which is negligible compared with $P(R)$. Its contribution from $L\ge n^{1+2\delta}$ is at most the already proved upper bound

\[
Q(n^{1+2\delta})\le n^{-1/2-\delta+o(1)},
\]

also negligible. Consequently the remaining polygons, with

\[
n\le L<n^{1+2\delta},\qquad D\ge R,
\]

carry length-weighted mass at least $n^{-(1+\delta)/2+o(1)}$. This is a lower bound for $Q(n)$. Dividing it by $n^{1+2\delta}$ gives

\[
F(n)\ge n^{-3/2-5\delta/2+o(1)}.
\]

Letting $\delta$ decrease to zero proves both claimed tail exponents.

Now condition on $L\ge n$. It suffices to prove the claim for $0<\epsilon<3/4$, since any larger tolerance follows from a smaller one. Fix such an $\epsilon$. For the unrooted law, the probability of an excessively large diameter is bounded by

\[
\frac{A(n^{3/4+\epsilon})}{F(n)}
       \le n^{-2\epsilon+o(1)}.
\]

The probability of an excessively small diameter is bounded by

\[
\frac{n^{-2}M_2(n^{3/4-\epsilon})}{F(n)}
       \le n^{-2\epsilon/3+o(1)}.
\]

Both vanish. Under the length-weighted law the corresponding bounds are

\[
\frac{P(n^{3/4+\epsilon})}{Q(n)}
       \le n^{-2\epsilon/3+o(1)},\qquad
 \frac{n^{-1}M_2(n^{3/4-\epsilon})}{Q(n)}
       \le n^{-2\epsilon/3+o(1)}.
\]

The second inequality uses $L\le L^2/n$ on the conditioning event. Thus $D=n^{3/4+o(1)}$ in probability under either law.

Finally condition the length-weighted law on $D\ge R$. Lengths below $R^{4/3-\epsilon}$ have conditional mass at most

\[
\frac{R^{4/3-\epsilon}A(R)}{P(R)}
       \le R^{-\epsilon+o(1)}.
\]

Lengths above $R^{4/3+\epsilon}$ have conditional mass at most

\[
\frac{Q(R^{4/3+\epsilon})}{P(R)}
       \le R^{-\epsilon/2+o(1)}.
\]

This proves $L=R^{4/3+o(1)}$ in probability.

The same conclusions hold with the length cutoff replaced by a power-separated window $[n,n^{1+\delta})$ for any fixed $\delta>0$, because the mass above its upper endpoint is negligible compared with the mass above $n$. Nothing in this proof estimates the mass at one exact length. That additional question requires a local estimate or a separate conditioning argument.
\end{proof}

\subsection{Open paths with fixed nearby endpoints}\label{sec:open-cutoffs}

For plane SAWs between the two vertices of a prescribed nearest-neighbor
bond, adding the bond closes every walk except the one-edge walk into
a polygon. This is
an exact correspondence with polygons containing the fixed edge, up to
the constant length and weight adjustment. For restrictions invariant
under honeycomb symmetries, the fixed-edge polygon mass is one third of
the length-weighted translation-class sum: symmetry equates the three
edge-type totals, and translation within a type enumerates the edges
of each class. Thus these open walks have the typical diameter law under
a large length cutoff and the typical length law under a large diameter
cutoff proved above for length-weighted polygons.

The half-plane walks of Corollary~\ref{cor:boundary-neighbor-mean} use a
different correspondence. Their bounded-change cap injection transfers
\eqref{c4:Geom-3} as an upper bound, in addition to the two diameter tails
already obtained there. Applying the same positive-tail argument gives
length-tail exponent $-3/2+o(1)$ under their critical edge weights,
and typical diameter $n^{3/4+o(1)}$ conditional on length at least $n$.
This is a length-cutoff assertion; it does not assert typical length
conditional only on large diameter under the unrooted weight.
